\chapter{Geometry \& Topology}

This section provides prerequisite knowledge, fundamental definitions, core theorems, and illustrative examples for Geometry and Topology, catering to beginner graduate students or advanced undergraduates.

\section{Algebraic Topology}
\textbf{Prerequisite Knowledge:} Set theory, point-set topology (open/closed sets, compactness, connectedness), and group theory.

\textbf{Definitions:}
\begin{itemize}
    \item \textbf{Fundamental Group ($\pi_1(X, x_0)$):} The group of equivalence classes under homotopy of loops contained in a topological space $X$ based at a point $x_0$.
    \item \textbf{Covering Space:} A continuous surjective map $p: E \to B$ is a covering space if every point $x \in B$ has an open neighborhood $U$ such that $p^{-1}(U)$ is a disjoint union of open sets in $E$, each mapped homeomorphically onto $U$ by $p$.
    \item \textbf{Singular Homology ($H_n(X)$):} A sequence of abelian groups associated to a topological space $X$, constructed from continuous maps of standard $n$-simplices into $X$.
    \item \textbf{CW Complex:} A topological space constructed by adjoining cells of increasing dimensions. It provides a highly computable framework for homology (Cellular Homology).
\end{itemize}

\textbf{Theorems:}
\begin{itemize}
    \item \textbf{Brouwer Fixed-Point Theorem:} Every continuous function from a closed $n$-ball $D^n$ to itself has at least one fixed point.
    \item \textbf{Mayer-Vietoris Sequence:} A long exact sequence in homology that relates the homology groups of a space $X = U \cup V$ to the homology groups of $U$, $V$, and their intersection $U \cap V$.
    \item \textbf{Hopf Fibration ($S^1 \hookrightarrow S^3 \to S^2$):} A non-trivial fiber bundle of the 3-sphere over the 2-sphere with fiber $S^1$. It famously shows that higher homotopy groups of spheres, such as $\pi_3(S^2)$, are non-trivial.
\end{itemize}

\textbf{Examples:}
\begin{itemize}
    \item \textbf{Example:} The fundamental group of the circle is $\pi_1(S^1) \cong \mathbb{Z}$. The integer corresponds to the winding number of a loop around the circle.
    \item \textbf{Example:} For any space $X$, the reduced homology of its suspension $\Sigma X$ satisfies $\widetilde{H}_{n}(\Sigma X) \cong \widetilde{H}_{n-1}(X)$.
\end{itemize}

\section{Differential Topology}
\textbf{Prerequisite Knowledge:} Advanced calculus (multivariable differentiation, implicit function theorem), point-set topology.

\textbf{Definitions:}
\begin{itemize}
    \item \textbf{Smooth Manifold:} A topological manifold equipped with a smooth atlas, allowing the calculus of smooth functions.
    \item \textbf{Vector Bundle:} A family of vector spaces parameterized continuously by a topological space $X$. The Tangent Bundle $TM$ of a smooth manifold $M$ is the primary example.
    \item \textbf{Characteristic Classes:} Cohomology classes associated with vector bundles that measure how much a bundle deviates from being trivial (e.g., Stiefel-Whitney classes for real bundles, Euler class for oriented bundles, Chern classes for complex bundles).
    \item \textbf{Cobordism:} Two compact $n$-dimensional manifolds $M$ and $N$ are \textit{cobordant} if their disjoint union is the boundary of a compact $(n+1)$-dimensional manifold $W$.
\end{itemize}

\textbf{Theorems:}
\begin{itemize}
    \item \textbf{Poincar\'e-Hopf Index Theorem:} If $M$ is a compact differentiable manifold and $V$ is a vector field with isolated zeros, the sum of the indices at the zeros equals the Euler characteristic $\chi(M)$.
    \item \textbf{Pontryagin-Thom Construction:} Establishes a bijection between the set of framed cobordism classes of submanifolds and the stable homotopy groups of spheres.
\end{itemize}

\textbf{Examples:}
\begin{itemize}
    \item \textbf{Example:} By the Hairy Ball Theorem (an application of Poincar\'e-Hopf), any continuous tangent vector field on $S^2$ must vanish somewhere.
    \item \textbf{Example:} The Real Projective Plane $\mathbb{RP}^2$ is unorientable, and thus has a non-zero first Stiefel-Whitney class $w_1(\mathbb{RP}^2) \neq 0$.
\end{itemize}

\section{Differential Geometry}
\textbf{Prerequisite Knowledge:} Linear algebra, multivariable calculus, ODEs.

\textbf{Definitions:}
\begin{itemize}
    \item \textbf{Riemannian Metric:} A smooth family of inner products on the tangent spaces of a differentiable manifold.
    \item \textbf{Levi-Civita Connection:} The unique affine connection on the tangent bundle of a Riemannian manifold that is symmetric (torsion-free) and compatible with the metric.
    \item \textbf{Curvature:} Quantities describing the deviation of a geometric object from being flat.
    \item \textbf{Mean Curvature:} The average of the principal curvatures of a surface. A surface with zero mean curvature everywhere is a \textit{minimal surface}.
\end{itemize}

\textbf{Theorems:}
\begin{itemize}
    \item \textbf{Fundamental Theorem of Riemannian Geometry:} On any Riemannian manifold, there is a unique Levi-Civita connection.
    \item \textbf{Gauss-Bonnet Theorem:} Links the local geometry of a surface to its global topology: $\int_{M} K \,dA = 2\pi \chi(M)$, where $K$ is the Gaussian curvature.
    \item \textbf{Myers's Theorem:} If a complete Riemannian manifold has Ricci curvature bounded below by a positive constant, its diameter is finite, and its fundamental group is finite.
\end{itemize}

\textbf{Examples:}
\begin{itemize}
    \item \textbf{Example:} The sphere $S^n$ has constant positive sectional curvature; Euclidean space $\mathbb{R}^n$ is flat; Hyperbolic space $\mathbb{H}^n$ has constant negative curvature.
\end{itemize}

\section{Lie Groups and Transformation Groups}
\textbf{Prerequisite Knowledge:} Linear algebra, manifold theory.

\textbf{Definitions:}
\begin{itemize}
    \item \textbf{Lie Group:} A smooth manifold $G$ that is also a group such that the multiplication and inversion operations are smooth maps.
    \item \textbf{Lie Algebra:} The tangent space to a Lie group at the identity element, equipped with a non-associative, alternating bilinear operation called the Lie bracket $[X, Y]$.
    \item \textbf{Group Action:} A Lie group $G$ acting smoothly on a manifold $M$.
\end{itemize}

\textbf{Theorems:}
\begin{itemize}
    \item \textbf{Topology of SO(n):} The fundamental group of the special orthogonal group $SO(n)$ for $n \ge 3$ is $\pi_1(SO(n)) \cong \mathbb{Z}/2\mathbb{Z}$.
    \item \textbf{Cartan-Dieudonn\'e Theorem:} Every orthogonal transformation in an $n$-dimensional symmetric bilinear space can be described as the composition of at most $n$ reflections.
\end{itemize}

\textbf{Examples:}
\begin{itemize}
    \item \textbf{Example:} $SO(3) \cong \mathbb{RP}^3$, showing that $SO(3)$ is doubly covered by $SU(2) \cong S^3$.
\end{itemize}

\section{Global Analysis and Geometric PDE}
\textbf{Prerequisite Knowledge:} Functional analysis, PDE, manifold theory.

\textbf{Definitions:}
\begin{itemize}
    \item \textbf{Hodge Star Operator:} A linear map on the exterior algebra of a finite-dimensional oriented inner product space, mapping $k$-forms to $(n-k)$-forms.
    \item \textbf{Harmonic Forms:} Differential $k$-forms $\omega$ that are in the kernel of the Laplace-Beltrami operator.
    \item \textbf{Elliptic Operator:} Differential operators that generalize the Laplacian, associated with "smoothing" properties.
\end{itemize}

\textbf{Theorems:}
\begin{itemize}
    \item \textbf{Hodge Decomposition Theorem:} Any differential $k$-form on a compact Riemannian manifold can be uniquely decomposed into the sum of an exact form, a co-exact form, and a harmonic form.
    \item \textbf{Bochner Technique:} A method of proving that certain cohomology groups vanish by using the Bochner formula.
\end{itemize}

\textbf{Examples:}
\begin{itemize}
    \item \textbf{Example:} The signature of 4-manifolds is an important topological invariant derived from the intersection form on the middle cohomology group $H^2(M; \mathbb{R})$.
\end{itemize}
